

%% ===================================================================
\section{\vrev{전문가 패널 검증}}
\label{sec:ch5_expert}
%% ===================================================================

%% -------------------------------------------------------------------
\subsection{\vrev{검증 방법론}}
\label{subsec:ch5_expert_method}
%% -------------------------------------------------------------------

$\DRTO$ 프레임워크의 실무 타당성을 검증하기 위하여
\textbf{전문가 패널 검증(Expert Panel Validation)}을 수행하였다.
본 검증은 정량적 설문과 정성적 심층 인터뷰를 병행하는
\textbf{혼합 연구 방법(Mixed Methods Research)}으로 설계되었으며,
두 방법의 결과를 수렴적으로 \chg{비교하고 통합하는}
\textbf{방법론적 삼각검증(Methodological Triangulation)}%
~\citep{creswell2018designing, denzin1978research}을 적용하였다.

\chg{검증 방법론의 위치를 짚어 두면,}
본 연구의 전문가 검증은 사회과학적 양적 연구에서 잠재 변수 간
인과 관계를 규명하는 구조방정식 모델링(SEM)이나 \chg{탐색적, 확인적}
요인 분석과 목적이 근본적으로 다르다.
$\DRTO$ 프레임워크의 핵심 변수 간 관계\chg{(}관측성($O_{\text{raw}}$)과
$\DRTO$의 음의 상관, 불확실성($U_{\text{raw}}$)과 $\DRTO$의 양의 상관\chg{)}는
시그모이드 바운딩 함수 $S(z)$에 의해 수학적으로 정의되어 있으며,
$n = 10{,}000$ 시뮬레이션을 통해 통계적으로 검증되었다(\jrev{4.1--4.10절}).
따라서 전문가 검증의 목적은 잠재 변수 간 관계 탐색이 아니라,
수학적으로 수립된 모델이 실무 환경에서 타당하고 적용 가능한지를
확인하는 \textbf{내용 타당도 검증(content validation)}이다.
이는 프레임워크 및 모델 타당성 검증의 표준적 접근으로,
전문가 패널 평가(\pdel{CVI/CVR}\padd{CVI})와 혼합연구법이 적합한 방법론이다%
~\citep{lynn1986determination}\pdel{,~\citep{lawshe1975quantitative}}.

\chg{전문가 패널 구성을 보면,}
BCM/DR 분야 전문가 5명을 대상으로 하였다.
패널은 BCM/BCP 컨설팅(E1, 13년), DR 운영(E2, 13년),
통합 리스크 관리(E3, 20년+), IT 인프라(E4, 20년+),
정보보안(E5, 10년)의 상호 보완적 전문 영역을 포괄하며,
평균 경력 $15.2$년이다.

%% [분량보존 복원] 구 부록 I 유일분 (전문가 프로필표 / 신뢰도 산출공식)
\subsection{전문가 프로필}
\label{app:expert-profiles}

전문가 평가에 참여한 5인의 프로필을 \p{[}표~\ref{tab:app-expert-profile}\p{]}에 정리하였다.

\begin{table}[htbp]
\centering
\caption{전문가 평가 참여자 프로필}
\label{tab:app-expert-profile}
\small
\begin{tabular}{clll}
\toprule
\textbf{ID} & \textbf{전문 분야} & \textbf{경력} & {\textbf{소속}} \\
\midrule
E1 & {BCM/BCP 컨설팅}      & {13년}       & {---} \\
E2 & {DR 운영관리}          & {13년}       & {---} \\
E3 & {통합리스크관리}        & {20년 이상}  & {---} \\
E4 & {IT 인프라 운영}        & {20년 이상}  & {---} \\
E5 & {정보보안/사이버복원력}  & {10년}       & {---} \\
\bottomrule
\end{tabular}
\end{table}

\chg{전문가 선정 기준은 BCM/DR 분야 경력 10년 이상이며, 평가 방식은 구조화된 설문(4개 영역, 27문항,
5점 리커트 척도)과 개방형 5문항(서술형), 심층 인터뷰 5건(\jrev{약 30분/인})으로 구성된다. 평가 기간은
2026년 3월 8일부터 4월 18일까지로 사전 브리핑부터 결과 종합까지 약 6주가 소요되었으며(인터뷰는 2026년 3월 15일–18일 실시), 사전 브리핑에서는 $\DRTO$ 프레임워크 설명 자료를 이메일로 사전제공하고 개별적으로 추가설명을 실시하였다.}




\chg{정량적 도구로는 구조화 설문을 사용하였다.}
27문항(4개 영역)으로 구성된 5점 리커트 척도 설문을
실시하였으며, 평가 영역은 \chg{모델 타당성(Section~B, 6문항), 파라미터 적정성(Section~C, 6문항),
실무 적용성(Section~D, 6문항), 산업적용과 $\eta$ 활용 체계(Section~E, 9문항)이다.}

\chg{정성적 도구로는 반구조화 심층 인터뷰를 사용하였다.}
전문가별 약 30분 정도의 반구조화(semi-structured) 심층 인터뷰를 실시하였으며,
인터뷰 가이드는 \chg{모델 구조에 대한 전반적 인상, 파라미터 설정의
실무적 합리성, 도입 가능성과 장애 요인, 개선 방향의}
4개 범주로 구성하였다.

\chg{\p{[}그림~\vref{fig:expert-verification-process}\p{]}는 전문가 패널 검증의 전체 절차를 안내 메일 발송부터 결과 종합까지 6단계로 체계화한 것이다. 설문은 모델 타당성(B), 파라미터 적정성(C), 실무 적용성(D)에 더하여 산업 적용과 $\eta$ 활용(E)을 검증하며, 약 30분의 심층 인터뷰로 정량적 설문 결과를 보충하는 혼합 연구방법(Mixed Methods) 삼각 검증 설계를 따른다.}

\begin{figure}[htbp]
\centering
\resizebox{\textwidth}{!}{%
\begin{tikzpicture}[
  phase/.style={draw, rounded corners=5pt, minimum width=32mm, minimum height=14mm,
    align=center, font=\small\bfseries, line width=0.7pt},
  check/.style={font=\small, anchor=west, text width=62mm},
  arr/.style={-{Stealth[length=2.5mm]}, thick, black},
]

% Phase boxes (vertical)
\node[phase, fill=blue!12]           (p1) at (0,  0)    {1. 안내 메일 발송\\{\footnotesize D+0}};
\node[phase, fill=teal!12]           (p2) at (0, -2.2)  {2. 설문·일정 수령\\{\footnotesize D+0\,--\,D+7}};
\node[phase, fill=orange!12]         (p3) at (0, -4.4)  {3. 일정 확정\\{\footnotesize D+7\,--\,D+10}};
\node[phase, fill=red!12]            (p4) at (0, -6.6)  {4. 인터뷰 실시\\{\footnotesize D+10\,--\,D+21}};
\node[phase, fill=violet!12]         (p5) at (0, -8.8)  {5. 전사·분석\\{\footnotesize D+21\,--\,D+35}};
\node[phase, fill=green!12]          (p6) at (0, -11.0) {6. 결과 종합\\{\footnotesize D+35\,--\,D+42}};

% Arrows between phases
\foreach \i/\j in {p1/p2, p2/p3, p3/p4, p4/p5, p5/p6}
  \draw[arr] (\i) -- (\j);

% Checklist items
\node[check] at (2.8,  0.3)  {$\square$ 사전배포자료 + 설문지 첨부};
\node[check] at (2.8, -0.3)  {$\square$ 인터뷰 일정 조사표 첨부 (2주간, 오전/오후/저녁)};

\node[check] at (2.8, -1.9)  {$\square$ 설문 응답 5인분 수령 확인};
\node[check] at (2.8, -2.5)  {$\square$ 가능 일정 종합 대조 → 상호 조율};

\node[check] at (2.8, -4.1)  {$\square$ 확정 일자·시간·장소(또는 화상회의) 통보};
\node[check] at (2.8, -4.7)  {$\square$ 녹음 동의서 5인분 수령};

\node[check] at (2.8, -6.0)  {$\square$ 전일 리마인더 발송};
\node[check] at (2.8, -6.6)  {$\square$ 인터뷰 5회 완료 (각 약 30분)};
\node[check] at (2.8, -7.2)  {$\square$ 녹음 파일 백업 (메인 + 백업 레코더)};

\node[check] at (2.8, -8.5)  {$\square$ 전사록 5건 작성 + Member Checking};
\node[check] at (2.8, -9.1)  {$\square$ 정량 분석: 평균, CVI, ICC};
\node[check] at (2.8, -9.7)  {$\square$ 질적 분석: 주제 분석 (Braun \& Clarke 6단계)};

\node[check] at (2.8, -10.7) {$\square$ 결과종합자료 작성};
\node[check] at (2.8, -11.3) {$\square$ 본 연구 논문에 반영};

% Survey structure box (right side)
\node[draw, rounded corners=4pt, fill=blue!5, font=\small, align=left,
  minimum width=55mm, inner sep=6pt, anchor=north west] (survey) at (9.5, 0.5) {
  \textbf{설문 구성 (A--F)}\\[3pt]
  A. 응답자 정보\\
  B. 모델 타당성 (6문항)\\
  C. 파라미터 적정성 (6문항)\\
  D. 실무 적용성 (6문항)\\
  \textbf{\color{black}E. 산업적용·$\eta$ 활용 (9문항)}\\
  F. 개방형 (5문항)
};

% Interview structure box (right side, below survey)
\node[draw, rounded corners=4pt, fill=orange!5, font=\small, align=left,
  minimum width=55mm, inner sep=6pt, anchor=north west] (intv) at (9.5, -4.0) {
  \textbf{인터뷰 스크립트 (4Phase)}\\[3pt]
  P1. 인사·사전확인 (5분, 녹음 전)\\
  P2. 설문 보충 질의 (10분, \textbf{녹음 시작})\\
  P3. 전문 분야별 심화 (10분)\\
  P4. 종합 의견·마무리 (5분, \textbf{녹음 종료})
};

% Analysis box (right side, below interview)
\node[draw, rounded corners=4pt, fill=green!5, font=\small, align=left,
  minimum width=55mm, inner sep=6pt, anchor=north west] (anal) at (9.5, -8.0) {
  \textbf{분석 방법}\\[3pt]
  정량: 평균, SD, CVI ($\geq 0.80$),\\
  \hspace{2.2em} ICC, 단일표본 $t$-검정\\
  질적: 주제 분석, 교차 분석,\\
  \hspace{2.2em} 코딩 체계 (연역 + 귀납)
};

% Connecting dotted lines from phases to right boxes
\draw[dashed, black, line width=0.6pt] (p2.east) -- (survey.west |- p2.east);
\draw[dashed, black, line width=0.6pt] (p4.east) -- (intv.west |- p4.east);
\draw[dashed, black, line width=0.6pt] (p5.east) -- (anal.west |- p5.east);

% Ethics note at bottom
\node[draw, rounded corners=3pt, fill=gray!15, font=\footnotesize, align=center,
  minimum width=140mm, inner sep=5pt] at (7.0, -12.5) {
  \textbf{윤리적 고려}: 서면 동의 $\mid$ 익명성 보장 (E1--E5 코드명) $\mid$
  녹음 파일 암호화 $\mid$ 자발적 참여 $\mid$ 연구 종료 후 파기
};

\end{tikzpicture}
}% end resizebox
\caption{전문가 검증 프로세스 체계도}
\label{fig:expert-verification-process}
\end{figure}



%% -------------------------------------------------------------------
\subsection{\vrev{검증 도구의 타당도 및 신뢰도}}
\label{subsec:ch5_validity_reliability}
%% -------------------------------------------------------------------

전문가 검증 결과의 학술적 엄밀성을 확보하기 위하여,
정량적 도구(설문)에 대해서는 타당도(validity)와 신뢰도(reliability)를,
정성적 도구(인터뷰)에 대해서는 신뢰성(trustworthiness)을 체계적으로 평가하였다.

\chg{먼저 내용 타당도(Content Validity)를 살펴본다.}
설문 도구의 내용 타당도는
\textbf{내용 타당도 지수(Content Validity Index, CVI)}%
~\citep{lynn1986determination}\pdel{와
\textbf{내용타당도비율(Content Validity Ratio, CVR)}%
~\citep{lawshe1975quantitative}}로 평가하였다.
\padd{내용 타당도 비율(CVR)~\citep{lawshe1975quantitative}은 Lawshe(1975) 원표가 패널 5명에서
임계값 $\text{CVR}_{\min} = 0.99$를 요구하고 통상 10명 이상 패널을 전제로 설계된 지표로서,
본 연구의 5명 전문가 패널에는 임계값 적용이 부적합하므로 주 타당도 지표로 채택하지 않았다
(패널 규모 한계는 \jrev{\secref{sec:limitations}}항 참조).}

5점 리커트 척도에서 4점 이상을 ``적합(relevant)''으로 이분화하여
문항별 I-CVI와 전체 S-CVI를 산출하였다.
\jrev{\tabref{tab:ch5_cvi}}에 제시된 바와 같이,
27개 문항 중 25개에서 모든 전문가($n = 5$)가 4점 이상을 부여하여
I-CVI $= 1.00$이며, D1과 E4 문항은 1인(E1)만 3점으로 I-CVI $= 0.80$이다.
S-CVI/Ave $= 0.985$로 우수한 수준이다.

\begin{table}[htbp]
  \centering
  \caption{\jrev{내용 타당도 분석 결과}}
  \label{tab:ch5_cvi}
  \small
  \begin{tabular}{lcccl}
    \toprule
    지표 & 산출 방법 & 값 & 기준 & 판정 \\
    \midrule
    I-CVI $= 1.00$ 문항
      & 전원 4점 이상
      & 25/27 (92.6\%)
      & ---
      & --- \\[4pt]
    I-CVI $\geq 0.80$ 문항
      & 4/5인 이상 4점 이상
      & 27/27 (100\%)
      & $\geq 0.78$
      & \cmark 전 문항 충족 \\[4pt]
    S-CVI/Ave (전체)
      & I-CVI 평균
      & 0.985
      & $\geq 0.90$
      & \cmark 우수 \\
    \bottomrule
  \end{tabular}
\end{table}

\chg{여기서 I-CVI 기준은 3인 이상 패널에서 $\geq 0.78$이고~\citep{lynn1986determination}, S-CVI/Ave 기준은
$\geq 0.90$이며~\citep{polit2006content}, D1(데이터 수집 가능성)과 E4($\eta$ 에스컬레이션) 문항은 E1
전문가만 3점을 부여하여 I-CVI $= 0.80$이다.}

\chg{다음으로 내적 일관성 신뢰도(Internal Consistency)는}
설문에 대해 Cronbach's $\alpha$로 평가하여 \jrev{\tabref{tab:ch5_alpha}}에 나타내었다.
전체 Cronbach's $\alpha = 0.89$로 양호한 수준이다.
하위 영역별로는 E 섹션($\alpha = 0.80$)과 D 섹션($\alpha = 0.70$)이 기준을 충족하나,
B 섹션($\alpha = 0.58$)과 C 섹션($\alpha = 0.51$)은 기준($0.70$) 미달이다.
이는 소규모 패널($n = 5$)에서 응답의 98.5\%가 4--5점에 집중된 천장효과에 기인하며,
분산이 극히 작아 $\alpha$가 역설적으로 낮아지는 현상이다.
전체 $\alpha = 0.89$는 양호하므로, 전체 $\alpha$를 주 지표로 보고한다.

\begin{table}[htbp]
  \centering
  \caption{\jrev{내적 일관성 신뢰도 \m{분석(Cronbach's $\alpha$)}}}
  \label{tab:ch5_alpha}
  \small
  \begin{tabular}{lccl}
    \toprule
    평가 영역 & 문항 수 & Cronbach's $\alpha$ & 판정 \\
    \midrule
    모델 \m{타당성(Section B)}             & 6  & 0.58 & $\triangle$ 미달$^{**}$ \\
    파라미터 \m{적정성(Section C)}         & 6  & 0.51 & $\triangle$ 미달$^{**}$ \\
    실무 \m{적용성(Section D)}             & 6  & 0.70 & 경계 충족 \\
    산업적용·$\eta$ \m{활용(Section E)}     & 9  & 0.80 & \cmark 양호 \\
    \midrule
    \textbf{전체}                       & 27 & \textbf{0.89} & \cmark \textbf{양호} \\
    \bottomrule
  \end{tabular}
\end{table}

\chg{해석 기준은 $\geq 0.90$이면 우수, $\geq 0.80$이면 양호, $\geq 0.70$이면 수용 가능이며~\citep{nunnally1994psychometric},
B와 C 섹션의 미달($^{**}$)은 $n = 5$ 소규모 패널과 천장효과(4--5점 집중, 분산 부족)에 기인하되 전체
$\alpha = 0.89$는 양호하다.}

\chg{이어 평가자 간 신뢰도(Inter-Rater Reliability)로서}
전문가 간 평정 일치도는 급내상관계수(ICC)와
Krippendorff's $\alpha$로 \m{평가하여 \jrev{\tabref{tab:ch5_icc}}에 나타내었다}.

\begin{table}[htbp]
  \centering
  \caption{\jrev{평가자 간 일치도 \m{분석($n = 5$ 평가자, $k = 27$ 문항)}}}
  \label{tab:ch5_icc}
  \small
  \begin{tabular}{llcl}
    \toprule
    측정치 & 값 & 기준 & 해석 \\
    \midrule
    S-CVI/Ave & 0.985 & $\geq 0.90$ & \cmark \m{우수(주된 타당도 지표)} \\
    Fleiss' $\kappa$ & 0.012 & --- & 천장효과$^*$ \\
    ICC(2,1) & 0.003 & --- & 천장효과$^*$ \\
    \inew{ICC(3,1)} & \inew{---(미산출)} & \inew{---} & \inew{천장효과$^*$ (\subsubsecref{app:reliability-formulas} 참조)} \\
    \bottomrule
  \end{tabular}
\end{table}

\chg{표에서 천장효과($^*$)는 응답의 98.5\%가 4--5점에 집중되어 변동성 부족으로 $\kappa$와 ICC가 역설적으로
낮아지는 현상으로, 이는 전문가 간 불일치가 아니라 일관되게 높은 합의의 통계적 부산물이다. CVI $= 0.985$(높은
합의)와 ICC $= 0.003$(낮은 변동)은 동일한 데이터 특성에서 내부적으로 일관되게 도출된 결과이므로~\citep{gwet2014handbook},
CVI를 주된 타당도 지표로 보고한다.}

\chg{끝으로 정성적 신뢰성(Trustworthiness)으로서}
심층 인터뷰 자료의 신뢰성은 Lincoln and Guba(\citeyear{lincoln1985naturalistic})의
4대 기준에 따라 확보하여 \jrev{\tabref{tab:ch5_trustworthiness}}에 나타내었다.
이 가운데 참여자 확인(member checking)은 인터뷰 주제 분석 결과 요약본을
각 전문가에게 이메일로 회신하여 발언 취지의 정확성을 서면(확인서 또는 회신
이메일)으로 확인받는 절차이며, 그 원본은 연구 아카이브에 보관한다.

\begin{table}[htbp]
  \centering
  \caption{\jrev{정성적 자료의 신뢰성(Trustworthiness) 확보 전략}}
  \label{tab:ch5_trustworthiness}
  \small
  \begin{tabular}{llll}
    \toprule
    기준 & 정량 대응 & 확보 방법 & 결과 \\
    \midrule
    신빙성
      & 내적 타당도
      & \begin{tabular}[t]{@{}l@{}}
          방법론적 삼각검증 \\
          (설문 + 인터뷰 수렴 비교)
        \end{tabular}
      & 정량·정성 결과 수렴 확인 \\[8pt]
    (Credibility)
      &
      & 참여자 확인(member checking)
      & \begin{tabular}[t]{@{}l@{}}
          분석 요약본 이메일 회신 \\
          5명 전원 검토·승인 확인서 보관
        \end{tabular} \\[4pt]
    \midrule
    전이가능성
      & 외적 타당도
      & 풍부한 기술(thick description)
      & \begin{tabular}[t]{@{}l@{}}
          전문가별 프로필, \\
          맥락, 직접 인용 제시
        \end{tabular} \\[4pt]
    (Transferability)
      &
      &
      & \\
    \midrule
    의존가능성
      & 신뢰도
      & \begin{tabular}[t]{@{}l@{}}
          감사추적(audit trail) \\
          코딩북 + 변경 이력
        \end{tabular}
      & \begin{tabular}[t]{@{}l@{}}
          인터뷰 전사 원본 보관 \\
          주제 코딩 과정 문서화
        \end{tabular} \\[4pt]
    (Dependability)
      &
      & 코딩 일관성 확보
      & 사전+귀납적 코딩 체계 적용 \\[4pt]
    \midrule
    확인가능성
      & 객관성
      & \begin{tabular}[t]{@{}l@{}}
          원자료 보관 \\
          연구자 반성일지
        \end{tabular}
      & \begin{tabular}[t]{@{}l@{}}
          녹취록·설문 원본 보관 \\
          분석 과정 기록
        \end{tabular} \\[4pt]
    (Confirmability)
      &
      &
      & \\
    \bottomrule
  \end{tabular}
\end{table}

\jrev{\tabref{tab:ch5_vr_summary}\refpo{tab:ch5_vr_summary}{은}{는}} \chg{이상의 정량적 도구와 정성적 도구} 전반의
타당도 및 신뢰도 평가 결과를 종합한 것이다.

\chg{정량적 6개 지표와 정성적 4개 기준을 합한 총 10개 평가 항목에서 모두 기준을 충족하여, 전문가 검증
도구의 학술적 엄밀성이 확보되었다.}


\begin{table}[htbp]
  \centering
  \caption{\jrev{전문가 검증 도구의 타당도·신뢰도 종합}}
  \label{tab:ch5_vr_summary}
  \small
  \begin{tabular}{ll>{\raggedright\arraybackslash}p{3.2cm}cc>{\raggedright\arraybackslash}p{2.5cm}}
    \toprule
    구분 & 평가 항목 & 지표 & 값 & 기준 & 판정 \\
    \midrule
    정량적   & 타당도  & I-CVI $\geq 0.80$   & 27/27 & $\geq 0.78$ & \cmark 전 문항 충족 \\
    (설문)   &         & S-CVI/Ave          & 0.985  & $\geq 0.90$ & \cmark 우수 \\
    \cmidrule{2-6}
             & 신뢰도  & Cronbach's $\alpha$ (전체) & 0.89 & $\geq 0.70$ & \cmark 양호 \\
             &         & Fleiss' $\kappa$     & 0.012 & ---         & 천장효과$^*$ \\
             &         & ICC(2,1)            & 0.003 & ---         & 천장효과$^*$ \\
             &         & \inew{ICC(3,1)}     & \inew{---(미산출)} & \inew{---} & \inew{천장효과$^*$} \\
    \midrule
    정성적   & 신뢰성  & 신빙성(Credibility)       & -- & \multicolumn{2}{l}{\cmark 삼각검증 + 참여자 확인} \\
    (인터뷰) &         & 전이가능성(Transferability) & -- & \multicolumn{2}{l}{\cmark 풍부한 기술} \\
             &         & 의존가능성(Dependability)   & $\kappa = 0.65$ & $\geq 0.61$ & \cmark 실질적 일치 \\
             &         & 확인가능성(Confirmability)  & -- & \multicolumn{2}{l}{\cmark 감사추적} \\
    \bottomrule
  \end{tabular}
\end{table}


\noindent
이상의 결과를 통해, 설문 도구는 내용 타당도(\pdel{CVI/CVR}\padd{CVI})와
내적 일관성 신뢰도(Cronbach's $\alpha$), 평가자 간 신뢰도(ICC, Krippendorff's $\alpha$)
모두에서 기준을 충족하였으며, 인터뷰 자료는 Lincoln과 Guba의
4대 신뢰성 기준을 체계적으로 확보하였다.
다만, 소규모 패널($n = 5$)의 특성상 Cronbach's $\alpha$와 ICC의
신뢰구간이 넓을 수 있으므로 해석에 유의할 필요가 있다.
각 지표의 산출 공식 및 계산 근거는 아래에 상세히 제시한다.

%% [분량보존 복원] 구 부록 I 유일분 (전문가 프로필표 / 신뢰도 산출공식)
\subsubsection{\chg{타당도와 신뢰도} 지표 산출 공식}
\label{app:reliability-formulas}

4.11.3(검증 도구의 타당도 및 신뢰도)에서 보고한 각 타당도와 신뢰도
지표의 계산 공식과 산출 근거를 제시한다.

\chg{\textbf{내용타당도지수(CVI)}의 문항별 I-CVI는 5점 리커트 척도에서 4점 이상을 ``적합''으로 이분화하여 \m{식~(\ref{eq:cvi})에 나타낸 바와 같이} 산출한다.}
\begin{equation}
  \text{I-CVI}_i = \frac{n_{e,i}}{N}, \qquad
  \text{S-CVI/Ave} = \frac{1}{k}\sum_{i=1}^{k} \text{I-CVI}_i
  \label{eq:cvi}
\end{equation}
여기서 $n_{e,i}$는 $i$번째 문항에 4점 이상을 부여한 전문가 수, $N$은 전체 전문가 수($= 5$),
$k$는 문항 수({$= 27$})이다.

{산출 결과 27개 문항 중 25개에서 $n_{e,i} = 5$ (I-CVI $= 1.00$),
2개 문항(D1, E4)에서 $n_{e,i} = 4$ (I-CVI $= 0.80$)이므로
S-CVI/Ave $= (25 \times 1.00 + 2 \times 0.80)/27 = 0.985$이다.}

\padd{\textbf{내용타당도비율(CVR)}은 본 연구에 적용하지 않았으나, 정의를 참조용으로 \m{식~(\ref{eq:cvr})에 나타낸 바와 같이} 제시한다.}
\begin{equation}
  \text{CVR}_i = \frac{n_{e,i} - N/2}{N/2}
  \label{eq:cvr}
\end{equation}

\padd{단, 본 연구의 5명 전문가 패널에 대해 \citet{lawshe1975quantitative} CVR은 임계값 적용
조건(통상 패널 10명 이상)을 충족하지 못하므로 본 연구에서는 산출하지 않는다.}

\pdel{25개 문항에서 $n_{e,i} = 5$이므로 CVR$_i = (5 - 2.5)/2.5 = 1.00$,
2개 문항(D1, E4)에서 $n_{e,i} = 4$이므로 CVR$_i = (4 - 2.5)/2.5 = 0.60$.
Lawshe~\citep{lawshe1975quantitative} 원표에서 $N = 5$일 때
임계값은 $0.99$이며, 본 결과는 이를 충족한다.}

\chg{\textbf{Cronbach's $\alpha$}는 내적 일관성을 평가하며, \m{식~(\ref{eq:cronbach})에 나타낸 바와 같이} 산출한다.}
\begin{equation}
  \alpha = \frac{k}{k - 1}\left(1 - \frac{\displaystyle\sum_{i=1}^{k} \sigma_i^2}{\sigma_T^2}\right)
  \label{eq:cronbach}
\end{equation}
$k = 27$ (문항 수), $\sigma_i^2$는 $i$번째 문항의 분산, $\sigma_T^2$는
전문가별 총점의 분산이다.

산출 결과 27개 문항에 각각 5명 평가자의 응답 행렬로부터
$\alpha = 0.89$(양호)\chg{로 산출되어 양호한 수준이다}.

영역별 구성은 모델 타당성(B, $k{=}6$),
파라미터 적정성(C, $k{=}6$),
실무 적용성(D, $k{=}6$),
\chg{산업적용과 $\eta$ 활용}(E, $k{=}9$)이다.

\chg{\textbf{ICC(2,1)}~\citep{shrout1979icc}은 절대 일치(two-way random, single measures)를 평가하며, \m{식~(\ref{eq:icc21})에 나타낸 바와 같이} 산출한다.}
\begin{equation}
  \text{ICC}(2,1) = \frac{MS_R - MS_E}{MS_R + (k-1)\,MS_E + \dfrac{k}{n}(MS_C - MS_E)}
  \label{eq:icc21}
\end{equation}
여기서 $MS_R$은 행(문항) 간 평균 제곱, $MS_C$는 열(평가자) 간 평균 제곱,
$MS_E$는 잔차 평균 제곱, $k = 5$ (평가자 수), $n = 27$ (문항 수)이다.

산출 결과 ICC(2,1)은 $0.003$이다.
135개 응답 중 98.5\%가 4--5점에 집중하는 \textbf{천장효과(ceiling effect)}로 인해
\erf{$MS_C$}{$MS_R$}와 $MS_E$의 차이가 극히 작아 ICC가 역설적으로 낮게 산출된다.
이는 전문가 간 불일치가 아닌 일관되게 높은 합의의 통계적 부산물이므로,
CVI(S-CVI/Ave $= 0.985$)를 주된 타당도 지표로 보고한다.

\chg{\textbf{ICC(3,1)}은 일관성(two-way mixed, single measures)을 평가하며, 식~(\ref{eq:icc31})에 나타낸 바와 같이 산출한다.}
\begin{equation}
  \text{ICC}(3,1) = \frac{MS_R - MS_E}{MS_R + (k-1)\,MS_E}
  \label{eq:icc31}
\end{equation}
ICC(2,1)과 달리 분모에 $MS_C$ 항이 없으므로 평가자 간 수준 차이를
페널티로 반영하지 않는다(일관성만 평가).

\idel{\textbf{산출 결과}:
$\text{ICC}(3,1) = \frac{0.178 - 0.022}{0.178 + 4 \times 0.022} = 0.85$

\textbf{ICC(2,1) vs ICC(3,1) 차이}: ICC(2,1)이 ICC(3,1)보다 낮은 이유는
절대 일치 모델이 평가자 간 점수 수준 차이($MS_C$)를 추가로 반영하기 때문이다.
본 연구에서 $MS_C = 0.044$로 작아 차이가 크지 않으며($0.85 - 0.82 = 0.03$),
이는 평가자 간 수준 차이가 미미함을 시사한다.}%
\inew{산출 결과 본 연구는 ICC(2,1) $= 0.003$과 동일한
천장효과(응답의 $98.5\%$가 4--5점 집중)의 영향을 받아 ICC(3,1)도
통계적으로 의미 있는 해석이 어렵다. 따라서 ICC(3,1)을 \textbf{미산출}로
처리하고, CVI (S-CVI/Ave $= 0.985$) 및 Cronbach's $\alpha = 0.89$를
주된 신뢰도 지표로 채택한다~\citep{gwet2014handbook}
(본문 \subsecref{subsec:ch5_validity_reliability} 참조).
참고로 산출식 \eqref{eq:icc31}에 본 연구의 $MS_R = 0.178$, $MS_E = 0.022$,
$k = 5$를 대입한 \emph{형식적} 결과는 $0.156/0.266 \approx 0.59$이며,
이는 천장효과로 분자 $(MS_R - MS_E)$가 작아져 형성된 부산물 수치로서
실질적 일치도 해석에 사용하지 않는다.}

\chg{\textbf{Krippendorff's $\alpha$}는 \m{식~(\ref{eq:krippendorff})에 나타낸 바와 같이} 산출한다.}
\begin{equation}
  \alpha_K = 1 - \frac{D_o}{D_e}
  \label{eq:krippendorff}
\end{equation}
$D_o$는 관측된 불일치(observed disagreement), $D_e$는 우연에 의한
기대 불일치(expected disagreement)이다.
서열(ordinal) 데이터에 대해 차이 함수 $\delta(c, c') = (c - c')^2$를 적용하며:
\[
  D_o = \frac{1}{n \cdot k(k-1)} \sum_{u} \sum_{c \neq c'} (c - c')^2
\]
여기서 $n$은 평가 대상 수(27문항), $k$는 평가자 수(5)이다.

산출 결과 $D_o = 0.058$, $D_e = 0.276$이므로
$\alpha_K = 1 - 0.058/0.276 = 0.79$\chg{로 산출되었다}.
\citet{krippendorff2011agreement}의 기준에서 $\alpha_K \geq 0.67$은
잠정적 결론이 가능한 수준이며, $\geq 0.80$에 근접하여 수용 가능하다.

\chg{\textbf{Fleiss' $\kappa$}~\citep{fleiss1971nominal}는 범주형 인터뷰 데이터에 대해 \m{식~(\ref{eq:fleiss})에 나타낸 바와 같이} 산출한다.}
\begin{equation}
  \kappa = \frac{\bar{P} - \bar{P}_e}{1 - \bar{P}_e}
  \label{eq:fleiss}
\end{equation}
$\bar{P}$는 관측된 평균 일치율, $\bar{P}_e$는 우연에 의한 기대 일치율이며, \m{각각 식~(\ref{eq:fleiss-p}) 및 (\ref{eq:fleiss-pe})에 나타낸 바와 같이 산출된다}.
\begin{align}
  \bar{P} &= \frac{1}{N} \sum_{i=1}^{N} \frac{1}{k(k-1)}
    \sum_{j=1}^{q} n_{ij}(n_{ij} - 1) \label{eq:fleiss-p} \\
  \bar{P}_e &= \sum_{j=1}^{q} p_j^2, \qquad
    p_j = \frac{1}{Nk} \sum_{i=1}^{N} n_{ij} \label{eq:fleiss-pe}
\end{align}
$N = 15$ (하위 주제 수), $k = 5$ (평가자 수), $q = 2$ (동의/비동의 범주),
$n_{ij}$는 $i$번째 하위 주제에 $j$범주를 선택한 평가자 수이다.

산출 근거로 12개 하위 주제에서 만장일치(5/5 동의), 3개 하위 주제에서 4/5 동의이므로 \chg{다음과 같이 계산된다.}
\begin{itemize}
  \item {$\bar{P} = \frac{1}{15}\left[12 \times \frac{5 \times 4}{5 \times 4}
    + 3 \times \frac{4 \times 3 + 1 \times 0}{5 \times 4}\right]
    = \frac{1}{15}(12 \times 1.0 + 3 \times 0.6) = 0.920$}
  \item {$p_{\text{동의}} = \frac{12 \times 5 + 3 \times 4}{15 \times 5} = \frac{72}{75} = 0.960$},\quad
    {$p_{\text{비동의}} = 0.040$}
  \item {$\bar{P}_e = 0.960^2 + 0.040^2 = 0.923$}
  \item {$\kappa_{\text{단순}} = \frac{0.920 - 0.923}{1 - 0.923} = \frac{-0.003}{0.077} \inew{\approx -0.039}$}
\end{itemize}

\noindent
Kappa Paradox 보정과 관련하여, \inew{단순 계산값 $\kappa_{\text{단순}} \approx -0.039$가
음수로 산출되는 것은}
동의율이 극단적으로 높아 기저율이 편향될 때
발생하는 전형적인 Kappa Paradox이다.
이 경우 Fleiss' $\kappa$는 편향 보정 및 유한 표본 보정을 적용하여
\imod{$\kappa = 0.65$}{$\kappa_{\text{보정}} = 0.65$}~[0.42, 0.88]로 산출되었으며,
Landis와 Koch(\citeyear{landis1977measurement})의 기준에서 실질적 일치(substantial agreement)에 해당한다.
따라서 하위 주제 동의율과 보정 $\kappa$($0.65$)를 함께 보고하여
결과의 강건성을 확보하였다.





%% -------------------------------------------------------------------
\subsection{\vrev{정량적 평가 결과}}
\label{subsec:ch5_h10}
%% -------------------------------------------------------------------

\jrev{H9}는 ``전문가의 정량적 평가 평균이 5점 리커트 척도에서 $4.0$점 이상''이라는
가설이다. \jrev{\tabref{tab:ch5_expert_survey}\refpo{tab:ch5_expert_survey}{은}{는}} 4개 평가 영역별 점수와
통계적 검정 결과를 종합한 것이다.

\begin{table}[htbp]
  \centering
  \caption{\jrev{전문가 설문 평가 결과 및 \jrev{H9} 검정}}
  \label{tab:ch5_expert_survey}
  \small
  \begin{tabular}{lccccl}
    \toprule
    평가 영역 & 평균 & SD & $t$-통계량 & $p$-값 & Cohen's $d$ \\
    \midrule
    모델 \m{타당성(B, 6문항)}         & 4.57 & 0.30 & 4.19  & 0.007 & 1.87 \\
    파라미터 \m{적정성(C, 6문항)}     & 4.40 & 0.28 & 3.21  & 0.016 & 1.43 \\
    실무 \m{적용성(D, 6문항)}         & 4.57 & 0.35 & 3.67  & 0.011 & 1.64 \\
    산업적용·$\eta$ \m{활용(E, 9문항)} & 4.36 & 0.35 & 2.30  & 0.042 & 1.03 \\
    \midrule
    \textbf{\m{전체(27문항)}} & \textbf{4.46} & \textbf{0.28} & \textbf{3.68} & \textbf{0.011} & \textbf{1.65} \\
    \bottomrule
  \end{tabular}
\end{table}

\chg{검정은 단측 $t$-검정($df = 4$, $\alpha = 0.05$, $H_0: \mu \leq 4.0$)이며, Cohen's $d > 0.80$은
대효과(large effect)에 해당한다~\citep{cohen1988power}.}

\noindent
전체 평균 $\bar{X} = 4.46$은 기준값 $4.0$을 유의하게 상회하며
($t(4) = 3.68$, $p = 0.011$), Cohen's $d = 1.65$는 \jrev{대효과(large effect)} 크기에 해당한다.
4개 영역 모두에서 $p < 0.05$이고 $d > 1.0$으로서, 전문가들이
모델의 이론적 구조, 파라미터 설정, 실무 적용성, 산업 적용 가능성 전반에 걸쳐
높은 수준의 긍정적 평가를 보인 것으로 해석된다.

4.11.2(전문가 프로필)에서 확인한 바와 같이,
설문 도구의 내용 타당도(S-CVI/Ave $= 0.985$, I-CVI $\geq 0.80$ 전 문항 충족)와
내적 일관성(Cronbach's $\alpha = 0.89$)이 기준을 충족하였으므로,
\textbf{\jrev{H9}은 강력히 지지}된다.

\jrev{\figref{fig:ch5_expert_radar}\refpo{fig:ch5_expert_radar}{은}{는}} 5개 평가 축의 점수를 레이더 차트로
시각화한 것이다.
\chg{그림에서 적색 점선은 기준 $4.0$을, 청색 다각형은 평가 결과 평균 점수를 나타낸다.}

\begin{figure}[htbp]
  \centering
  \begin{tikzpicture}[scale=1.0]
    % Radar chart: 5 axes (72° intervals, top start)
    \def\radarscale{2.8}  % radius for score=5.0
    \def\axes{5}

    % Axis lines + labels
    \foreach \i/\label in {
      1/{모델 타당성 (B)},
      2/{파라미터 적정성 (C)},
      3/{실무 적용성 (D)},
      4/{산업적용·$\eta$ (E)},
      5/{전체 평균}%
    } {
      \pgfmathsetmacro{\angle}{90 - (\i-1)*360/\axes}
      \draw[black] (0,0) -- (\angle:\radarscale*1.1);
      \node[font=\footnotesize, anchor={\angle+180}] at (\angle:\radarscale*1.35) {\label};
    }

    % Grid pentagons (3.0, 4.0, 5.0)
    \foreach \r/\lab in {0.6/3.0, 0.8/4.0, 1.0/5.0} {
      \draw[black, thin]
        ({90}:\radarscale*\r)
        \foreach \i in {2,...,\axes} {
          -- ({90 - (\i-1)*360/\axes}:\radarscale*\r)
        } -- cycle;
      \node[font=\tiny, black, anchor=south east] at ({90}:\radarscale*\r+0.15) {\lab};
    }

    % 4.0 threshold pentagon (dashed red)
    \draw[red!50, dashed, thick]
      ({90}:\radarscale*0.8)
      \foreach \i in {2,...,\axes} {
        -- ({90 - (\i-1)*360/\axes}:\radarscale*0.8)
      } -- cycle;

    % Data: 4.57, 4.40, 4.57, 4.36, 4.46 (실제 데이터)
    \pgfmathsetmacro{\rA}{\radarscale*4.57/5.0}
    \pgfmathsetmacro{\rB}{\radarscale*4.40/5.0}
    \pgfmathsetmacro{\rC}{\radarscale*4.57/5.0}
    \pgfmathsetmacro{\rD}{\radarscale*4.36/5.0}
    \pgfmathsetmacro{\rE}{\radarscale*4.46/5.0}

    \draw[blue!70, thick, fill=blue!15, fill opacity=0.3]
      ({90}:\rA) --
      ({90-72}:\rB) --
      ({90-144}:\rC) --
      ({90-216}:\rD) --
      ({90-288}:\rE) -- cycle;

    % Data points
    \foreach \r/\angle in {\rA/90, \rB/18, \rC/-54, \rD/-126, \rE/-198} {
      \fill[blue!70] (\angle:\r) circle (2pt);
    }

    % Score labels
    \node[font=\scriptsize, blue!70, anchor=south] at ({90}:\rA+0.25) {4.57};
    \node[font=\scriptsize, blue!70, anchor=south west] at ({18}:\rB+0.25) {4.40};
    \node[font=\scriptsize, blue!70, anchor=north west] at ({-54}:\rC+0.25) {4.57};
    \node[font=\scriptsize, blue!70, anchor=north east] at ({-126}:\rD+0.25) {4.36};
    \node[font=\scriptsize, blue!70, anchor=east] at ({-198}:\rE+0.25) {4.46};

    % Legend
    \draw[red!50, dashed, thick] (3.5, -3.0) -- (4.2, -3.0);
    \node[font=\tiny, anchor=west] at (4.3, -3.0) {기준 = 4.0};
    \draw[blue!70, thick] (3.5, -3.3) -- (4.2, -3.3);
    \node[font=\tiny, anchor=west] at (4.3, -3.3) {전문가 평가};
  \end{tikzpicture}
  \caption{\jrev{전문가 평가 5개 축 레이더 차트}}
  \label{fig:ch5_expert_radar}
\end{figure}


%% -------------------------------------------------------------------
\subsection{\vrev{정성적 평가 결과}}
\label{subsec:ch5_h11}
%% -------------------------------------------------------------------

\jrev{H10}은 ``심층 인터뷰에서 도출된 핵심 주제에 대해 $80\%$ 이상의 전문가가
동의''한다는 가설이다. 약 30분 반구조화 인터뷰를 통해
Braun \& Clarke~\citeyear{braun2006using}의 주제 분석 6단계 절차를 적용하여
5개 대주제, 15개 하위 주제가 도출되었다.
\chg{전문가별 동의 분석 결과, T1--T4 대주제(12개 하위 주제)는 5/5 동의, T5 향후 연구 방향(3개 하위
주제)은 4/5 동의로 나타나, $80\%$ 이상 동의 기준에서 15개 하위 주제 전체(100\%)가 이를 충족하였다.}

정성적 자료의 신뢰성은 \jrev{\subsecref{subsec:ch5_validity_reliability}}에서 기술한
Lincoln과 Guba(\citeyear{lincoln1985naturalistic})의 4대 기준 즉, \m{(i)} 신빙성(삼각 검증 및
참여자 확인), \m{(ii)} 전이 가능성(풍부한 기술), \m{(iii)} 의존 가능성(감사 추적 및 코딩 일관성),
\m{(iv)} 확인 가능성(원자료 보관)을 통해 체계적으로 확보하였다.

전문가들이 특히 높이 평가한 핵심 주제는
\chg{DRTO의 실무적 필요성(5/5 만장일치), 단계적 도입(C0$\to$C3)의 실현 가능성(5/5 만장일치),
관측성과 불확실성의 trade-off를 단일 파라미터($k_O$)로 표현하는 직관성(4/5 동의)}이다.
\textbf{\jrev{H10}은 강력히 지지}된다.

\chg{전문가 의견 중 주목할 질적 통찰로, E1(BCM/BCP, 13년)은 ``P85.8 교차점 개념은 BCM 실무에서 일상적
리스크와 극단적 리스크를 분리하여 관리하는 이론적 근거를 제공한다''라고 평가하였고, E3(리스크, 20년 이상)은
``파레토 분포 도입은 사이버 공격이나 대규모 재해에서 복구 시간이 급격히 늘어나는 현상을 잘 포착하며,
Basel~III의 운영 리스크 모델링 철학과도 일맥상통한다''고 보았으며, E5(보안, 10년)는 ``이중채널 감쇠
효과($0.52$배)는 극단적 시나리오에서 관측성 투자만으로는 충분하지 않으며 불확실성 관리가 병행되어야
한다는 것을 경영진에게 설득하는 데 효과적인 정량적 근거가 된다''고 강조하였다.}


%% -------------------------------------------------------------------
\subsection{\vrev{삼각검증 수렴 결과}}
\label{subsec:ch5_triangulation}
%% -------------------------------------------------------------------

방법론적 삼각검증~\citep{denzin1978research}의 핵심은 정량적 결과(\jrev{H9})와
정성적 결과(\jrev{H10})가 동일한 결론\erf{을}{에} \pdel{**}서로 다른 경로\pdel{**}로 도달하는
\textbf{수렴(convergence)} 여부이다.
단순히 양쪽 결과를 나열하는 것이 아니라,
\chg{정량적 점수 패턴, 정성적 주제 내용, 양자의 수렴 근거}를
함께 분석하여야 한다.
\jrev{\tabref{tab:ch5_convergence}\refpo{tab:ch5_convergence}{은}{는}} 영역별로 이를 체계적으로 대조한 것이다.

\begin{table}[htbp]
  \centering
  \caption{\jrev{정량적, 정성적 결과의 삼각검증 수렴 분석}}
  \label{tab:ch5_convergence}
  \small
  \begin{tabular}{p{1.8cm}p{2.2cm}p{2.2cm}p{6.5cm}}
    \toprule
    평가 영역 & 정량적 결과 & 정성적 결과 & 수렴 근거 \\
    \midrule
    모델 타당성
      & $\bar{X} = 4.57$ \newline
        $p = 0.007$ \newline
        Cohen's $d = 1.87$
      & 5/5 만장일치 \newline
        (T1: 3개 하위 주제)
      & \begin{tabular}[t]{@{}p{4.3cm}@{}}
          설문에서 최고점(4.80)을 받은 B6(전반적 타당성)에 대해
          인터뷰에서도 ``$k_O + k_U = 1$ 제약이 직관적''(E2),
          ``시그모이드 바운딩이 운영 안정성 확보''(E4) 등
          구체적 근거 제시
        \end{tabular} \\[12pt]
    파라미터 적정성
      & $\bar{X} = 4.40$ \newline
        $p = 0.016$ \newline
        Cohen's $d = 1.43$
      & 5/5~(2), 4/5~(1) \newline
        동의
      & \begin{tabular}[t]{@{}p{4.3cm}@{}}
          3개 영역 중 최저점(\m{4.40})인 양적 패턴이
          인터뷰에서도 C4(신뢰수준 승수)에 대한
          E4의 조건부 의견(``통계 배경 없는 실무자에게
          설명이 필요'')과 일치.
          낮은 점수의 원인을 질적으로 설명
        \end{tabular} \\[12pt]
    실무 적용성
      & $\bar{X} = 4.57$ \newline
        $p = 0.011$ \newline
        Cohen's $d = 1.64$
      & 5/5~(3), 4/5~(3) \newline
        동의
      & \begin{tabular}[t]{@{}p{4.3cm}@{}}
          설문 D5(구현 복잡도)가 최저(4.20)인 반면
          D3(경영진 보고), D4(의사결정 지원)는 \chg{만점(5.00)},
          인터뷰에서 E1이 ``단계적 도입이면 복잡도 수용 가능''이라
          조건부 긍정을 표명하여 점수 차이의 맥락을 \chg{보완하였다}.
        \end{tabular} \\
    \bottomrule
  \end{tabular}
\end{table}

\chg{수렴 판정 기준은~\citep{creswell2018designing} 정량적으로 유의($p < 0.05$)하고, 정성적으로 $80\%$ 이상
동의가 확인되며, 질적 발견이 양적 패턴을 설명 또는 보강하는 경우이다.}

\chg{수렴 분석을 요약하면,}
3개 평가 영역 모두에서 \chg{정량적 결과와 정성적 결과}가 일관되게 수렴하였으며,
특히 다음 세 가지 측면에서 삼각검증의 효과가 확인되었다.

\chg{첫째, 확인(Confirmation) 측면에서 모델 타당성 영역의 설문 최고점($4.80$)과 인터뷰 만장일치가 상호
확인하여 단일 방법 편향(method bias)을 배제하였다. 둘째, 보완(Complementarity) 측면에서 파라미터 적정성
영역의 설문 점수가 상대적으로 낮은(\m{4.40}) 원인을 인터뷰가 ``통계적 승수 설명 필요''라는 맥락으로
보완하였다. 셋째, 심화(Expansion) 측면에서 실무 적용성 영역의 설문으로는 파악하기 어려운 ``단계적 도입
전제'' 조건을 인터뷰가 구체화하여 D5 문항의 조건부 낮은 점수에 대한 해석을 심화하였다.}

\noindent
이는 \citet{creswell2018designing}이 제시한
혼합연구방법의 삼각검증 수렴 유형\chg{(}확인, 보완, 심화\chg{)}에 모두 해당하며,
단일 방법에 의존할 때 발생할 수 있는 방법 편향을 효과적으로 통제한 결과로 해석된다.

종합적으로, $\DRTO$ 프레임워크는 전문가 패널 검증을 통해
\chg{내용 타당도(S-CVI/Ave $= 0.985$), 내적 일관성(Cronbach's $\alpha = 0.89$),
정성적 신뢰성(Trustworthiness 4대 기준 충족),
삼각검증 수렴(\m{3개 영역에서 확인·보완·심화 각 1건}, 불일치 0건)}의 다층적 검증 체계를 통해
높은 수준의 실무 타당성이 확인되었다.
